\IfFileExists{stacks-project.cls}{%
\documentclass{stacks-project}
}{%
\documentclass{amsart}
}

% For dealing with references we use the comment environment
\usepackage{verbatim}
\newenvironment{reference}{\comment}{\endcomment}
%\newenvironment{reference}{}{}
\newenvironment{slogan}{\comment}{\endcomment}
\newenvironment{history}{\comment}{\endcomment}

% For commutative diagrams we use Xy-pic
\usepackage[all]{xy}

% We use 2cell for 2-commutative diagrams.
\xyoption{2cell}
\UseAllTwocells

% We use multicol for the list of chapters between chapters
\usepackage{multicol}

% This is generally recommended for better output
\usepackage{lmodern}
\usepackage[T1]{fontenc}

% For cross-file-references
\usepackage{xr-hyper}

% Package for hypertext links:
\usepackage{hyperref}

% For any local file, say "hello.tex" you want to link to please
% use \externaldocument[hello-]{hello}
\externaldocument[introduction-]{introduction}
\externaldocument[conventions-]{conventions}
\externaldocument[sets-]{sets}
\externaldocument[categories-]{categories}
\externaldocument[topology-]{topology}
\externaldocument[sheaves-]{sheaves}
\externaldocument[sites-]{sites}
\externaldocument[stacks-]{stacks}
\externaldocument[fields-]{fields}
\externaldocument[algebra-]{algebra}
\externaldocument[brauer-]{brauer}
\externaldocument[homology-]{homology}
\externaldocument[derived-]{derived}
\externaldocument[simplicial-]{simplicial}
\externaldocument[more-algebra-]{more-algebra}
\externaldocument[smoothing-]{smoothing}
\externaldocument[modules-]{modules}
\externaldocument[sites-modules-]{sites-modules}
\externaldocument[injectives-]{injectives}
\externaldocument[cohomology-]{cohomology}
\externaldocument[sites-cohomology-]{sites-cohomology}
\externaldocument[dga-]{dga}
\externaldocument[dpa-]{dpa}
\externaldocument[sdga-]{sdga}
\externaldocument[hypercovering-]{hypercovering}
\externaldocument[schemes-]{schemes}
\externaldocument[constructions-]{constructions}
\externaldocument[properties-]{properties}
\externaldocument[morphisms-]{morphisms}
\externaldocument[coherent-]{coherent}
\externaldocument[divisors-]{divisors}
\externaldocument[limits-]{limits}
\externaldocument[varieties-]{varieties}
\externaldocument[topologies-]{topologies}
\externaldocument[descent-]{descent}
\externaldocument[perfect-]{perfect}
\externaldocument[more-morphisms-]{more-morphisms}
\externaldocument[flat-]{flat}
\externaldocument[groupoids-]{groupoids}
\externaldocument[more-groupoids-]{more-groupoids}
\externaldocument[etale-]{etale}
\externaldocument[chow-]{chow}
\externaldocument[intersection-]{intersection}
\externaldocument[pic-]{pic}
\externaldocument[weil-]{weil}
\externaldocument[adequate-]{adequate}
\externaldocument[dualizing-]{dualizing}
\externaldocument[duality-]{duality}
\externaldocument[discriminant-]{discriminant}
\externaldocument[derham-]{derham}
\externaldocument[local-cohomology-]{local-cohomology}
\externaldocument[algebraization-]{algebraization}
\externaldocument[curves-]{curves}
\externaldocument[resolve-]{resolve}
\externaldocument[models-]{models}
\externaldocument[functors-]{functors}
\externaldocument[equiv-]{equiv}
\externaldocument[pione-]{pione}
\externaldocument[etale-cohomology-]{etale-cohomology}
\externaldocument[proetale-]{proetale}
\externaldocument[relative-cycles-]{relative-cycles}
\externaldocument[more-etale-]{more-etale}
\externaldocument[trace-]{trace}
\externaldocument[crystalline-]{crystalline}
\externaldocument[spaces-]{spaces}
\externaldocument[spaces-properties-]{spaces-properties}
\externaldocument[spaces-morphisms-]{spaces-morphisms}
\externaldocument[decent-spaces-]{decent-spaces}
\externaldocument[spaces-cohomology-]{spaces-cohomology}
\externaldocument[spaces-limits-]{spaces-limits}
\externaldocument[spaces-divisors-]{spaces-divisors}
\externaldocument[spaces-over-fields-]{spaces-over-fields}
\externaldocument[spaces-topologies-]{spaces-topologies}
\externaldocument[spaces-descent-]{spaces-descent}
\externaldocument[spaces-perfect-]{spaces-perfect}
\externaldocument[spaces-more-morphisms-]{spaces-more-morphisms}
\externaldocument[spaces-flat-]{spaces-flat}
\externaldocument[spaces-groupoids-]{spaces-groupoids}
\externaldocument[spaces-more-groupoids-]{spaces-more-groupoids}
\externaldocument[bootstrap-]{bootstrap}
\externaldocument[spaces-pushouts-]{spaces-pushouts}
\externaldocument[spaces-chow-]{spaces-chow}
\externaldocument[groupoids-quotients-]{groupoids-quotients}
\externaldocument[spaces-more-cohomology-]{spaces-more-cohomology}
\externaldocument[spaces-simplicial-]{spaces-simplicial}
\externaldocument[spaces-duality-]{spaces-duality}
\externaldocument[formal-spaces-]{formal-spaces}
\externaldocument[restricted-]{restricted}
\externaldocument[spaces-resolve-]{spaces-resolve}
\externaldocument[formal-defos-]{formal-defos}
\externaldocument[defos-]{defos}
\externaldocument[cotangent-]{cotangent}
\externaldocument[examples-defos-]{examples-defos}
\externaldocument[algebraic-]{algebraic}
\externaldocument[examples-stacks-]{examples-stacks}
\externaldocument[stacks-sheaves-]{stacks-sheaves}
\externaldocument[criteria-]{criteria}
\externaldocument[artin-]{artin}
\externaldocument[quot-]{quot}
\externaldocument[stacks-properties-]{stacks-properties}
\externaldocument[stacks-morphisms-]{stacks-morphisms}
\externaldocument[stacks-limits-]{stacks-limits}
\externaldocument[stacks-cohomology-]{stacks-cohomology}
\externaldocument[stacks-perfect-]{stacks-perfect}
\externaldocument[stacks-introduction-]{stacks-introduction}
\externaldocument[stacks-more-morphisms-]{stacks-more-morphisms}
\externaldocument[stacks-geometry-]{stacks-geometry}
\externaldocument[moduli-]{moduli}
\externaldocument[moduli-curves-]{moduli-curves}
\externaldocument[examples-]{examples}
\externaldocument[exercises-]{exercises}
\externaldocument[guide-]{guide}
\externaldocument[desirables-]{desirables}
\externaldocument[coding-]{coding}
\externaldocument[obsolete-]{obsolete}
\externaldocument[fdl-]{fdl}
\externaldocument[index-]{index}

% Theorem environments.
%
\theoremstyle{plain}
\newtheorem{theorem}[subsection]{Theorem}
\newtheorem{proposition}[subsection]{Proposition}
\newtheorem{lemma}[subsection]{Lemma}

\theoremstyle{definition}
\newtheorem{definition}[subsection]{Definition}
\newtheorem{example}[subsection]{Example}
\newtheorem{exercise}[subsection]{Exercise}
\newtheorem{situation}[subsection]{Situation}

\theoremstyle{remark}
\newtheorem{remark}[subsection]{Remark}
\newtheorem{remarks}[subsection]{Remarks}

\numberwithin{equation}{subsection}

% Macros
%
\def\lim{\mathop{\mathrm{lim}}\nolimits}
\def\colim{\mathop{\mathrm{colim}}\nolimits}
\def\Spec{\mathop{\mathrm{Spec}}}
\def\Hom{\mathop{\mathrm{Hom}}\nolimits}
\def\Ext{\mathop{\mathrm{Ext}}\nolimits}
\def\SheafHom{\mathop{\mathcal{H}\!\mathit{om}}\nolimits}
\def\SheafExt{\mathop{\mathcal{E}\!\mathit{xt}}\nolimits}
\def\Sch{\mathit{Sch}}
\def\Mor{\mathop{\mathrm{Mor}}\nolimits}
\def\Ob{\mathop{\mathrm{Ob}}\nolimits}
\def\Sh{\mathop{\mathit{Sh}}\nolimits}
\def\NL{\mathop{N\!L}\nolimits}
\def\CH{\mathop{\mathrm{CH}}\nolimits}
\def\proetale{{pro\text{-}\acute{e}tale}}
\def\etale{{\acute{e}tale}}
\def\QCoh{\mathit{QCoh}}
\def\Ker{\mathop{\mathrm{Ker}}}
\def\Im{\mathop{\mathrm{Im}}}
\def\Coker{\mathop{\mathrm{Coker}}}
\def\Coim{\mathop{\mathrm{Coim}}}

% Boxtimes
%
\DeclareMathSymbol{\boxtimes}{\mathbin}{AMSa}{"02}

%
% Macros for moduli stacks/spaces
%
\def\QCohstack{\mathcal{QC}\!\mathit{oh}}
\def\Cohstack{\mathcal{C}\!\mathit{oh}}
\def\Spacesstack{\mathcal{S}\!\mathit{paces}}
\def\Quotfunctor{\mathrm{Quot}}
\def\Hilbfunctor{\mathrm{Hilb}}
\def\Curvesstack{\mathcal{C}\!\mathit{urves}}
\def\Polarizedstack{\mathcal{P}\!\mathit{olarized}}
\def\Complexesstack{\mathcal{C}\!\mathit{omplexes}}
% \Pic is the operator that assigns to X its picard group, usage \Pic(X)
% \Picardstack_{X/B} denotes the Picard stack of X over B
% \Picardfunctor_{X/B} denotes the Picard functor of X over B
\def\Pic{\mathop{\mathrm{Pic}}\nolimits}
\def\Picardstack{\mathcal{P}\!\mathit{ic}}
\def\Picardfunctor{\mathrm{Pic}}
\def\Deformationcategory{\mathcal{D}\!\mathit{ef}}
\usepackage{cleveref}
\crefname{originalreference}{original reference}{original references}
\Crefname{originalreference}{Original reference}{Original references}
\makeatletter
\newlabel{algebra-lemma-compare-resolutions}{{00LS}{source}{}{sourcepointer-1}{}}
\newlabel{algebra-lemma-compare-resolutions@cref}{{[originalreference][0][]00LS}{[1][0][]source}{}{}{}}
\makeatother
\crefname{originalreference}{original reference}{original references}
\Crefname{originalreference}{Original reference}{Original references}
\makeatletter
\newlabel{algebra-lemma-grothendieck-regular-sequence-general}{{0470}{source}{}{sourcepointer-2}{}}
\newlabel{algebra-lemma-grothendieck-regular-sequence-general@cref}{{[originalreference][0][]0470}{[1][0][]source}{}{}{}}
\makeatother
\crefname{originalreference}{original reference}{original references}
\Crefname{originalreference}{Original reference}{Original references}
\makeatletter
\newlabel{more-algebra-definition-koszul-complex}{{0623}{source}{}{sourcepointer-3}{}}
\newlabel{more-algebra-definition-koszul-complex@cref}{{[originalreference][0][]0623}{[1][0][]source}{}{}{}}
\makeatother
\crefname{originalreference}{original reference}{original references}
\Crefname{originalreference}{Original reference}{Original references}
\makeatletter
\newlabel{more-algebra-lemma-lift-complex-stably-frees}{{0BCB}{source}{}{sourcepointer-4}{}}
\newlabel{more-algebra-lemma-lift-complex-stably-frees@cref}{{[originalreference][0][]0BCB}{[1][0][]source}{}{}{}}
\makeatother
\crefname{originalreference}{original reference}{original references}
\Crefname{originalreference}{Original reference}{Original references}
\makeatletter
\newlabel{more-algebra-lemma-regular-koszul-regular}{{062F}{source}{}{sourcepointer-5}{}}
\newlabel{more-algebra-lemma-regular-koszul-regular@cref}{{[originalreference][0][]062F}{[1][0][]source}{}{}{}}
\makeatother
\title{Koszul lifting along essentially finite presentation flat maps}
\author{The Stacks Project Authors}
\date{stacks-project:equiv.tex:local-snapshot-20261012; GNU FDL 1.2 or later}
\begin{document}
\maketitle
\noindent\textit{Allowed standard prerequisites (editorial).}\begin{itemize}
\item Stably free lifting. more-algebra-lemma-lift-complex-stably-frees (cross-chapter, named). Represents K by a bounded complex of free modules tensoring to the Koszul complex.
\item Regular sequences. algebra-lemma-grothendieck-regular-sequence-general (cross-chapter, named). Lifts form a regular sequence with flat quotient.
\item Resolution comparison. algebra-lemma-compare-resolutions (cross-chapter, named). Produces the dotted comparison arrow.
\end{itemize}\medskip
\ifdefined\hypertarget\hypertarget{sourcepointer-1}{}\fi
\noindent\textit{Source reference.} Original reference 00LS: focused-seeds/geometry/S9-eqv-deform-koszul/ref-tags.txt, source lines 5--5. Navigation pointer only: the counted proof cites this statement by \textbackslash{}ref as a standing input; the deduction the counted proof has to make is carried out inside the counted span, and no part of that deduction is delegated to this pointer.\par\medskip
\ifdefined\hypertarget\hypertarget{sourcepointer-2}{}\fi
\noindent\textit{Source reference.} Original reference 0470: focused-seeds/geometry/S9-eqv-deform-koszul/ref-tags.txt, source lines 6--6. Navigation pointer only: the counted proof cites this statement by \textbackslash{}ref as a standing input; the deduction the counted proof has to make is carried out inside the counted span, and no part of that deduction is delegated to this pointer.\par\medskip
\ifdefined\hypertarget\hypertarget{sourcepointer-3}{}\fi
\noindent\textit{Source reference.} Original reference 0623: focused-seeds/geometry/S9-eqv-deform-koszul/ref-tags.txt, source lines 7--7. Navigation pointer only: the counted proof cites this statement by \textbackslash{}ref as a standing input; the deduction the counted proof has to make is carried out inside the counted span, and no part of that deduction is delegated to this pointer.\par\medskip
\ifdefined\hypertarget\hypertarget{sourcepointer-4}{}\fi
\noindent\textit{Source reference.} Original reference 0BCB: focused-seeds/geometry/S9-eqv-deform-koszul/ref-tags.txt, source lines 8--8. Navigation pointer only: the counted proof cites this statement by \textbackslash{}ref as a standing input; the deduction the counted proof has to make is carried out inside the counted span, and no part of that deduction is delegated to this pointer.\par\medskip
\ifdefined\hypertarget\hypertarget{sourcepointer-5}{}\fi
\noindent\textit{Source reference.} Original reference 062F: focused-seeds/geometry/S9-eqv-deform-koszul/ref-tags.txt, source lines 9--9. Navigation pointer only: the counted proof cites this statement by \textbackslash{}ref as a standing input; the deduction the counted proof has to make is carried out inside the counted span, and no part of that deduction is delegated to this pointer.\par\medskip
\begin{lemma}
\label{lemma-deform-koszul}
Let $(R, \mathfrak m, \kappa) \to (A, \mathfrak n, \lambda)$
be a flat local ring homorphism of local rings
which is essentially of finite presentation.
Let $\overline{f}_1, \ldots, \overline{f}_r \in \mathfrak n/\mathfrak m A
\subset A/\mathfrak m A$ be a regular sequence. Let $K \in D(A)$. Assume
\begin{enumerate}
\item $K$ is perfect,
\item $K \otimes_A^\mathbf{L} A/\mathfrak m A$ is isomorphic in
$D(A/\mathfrak m A)$ to the
Koszul complex on $\overline{f}_1, \ldots, \overline{f}_r$.
\end{enumerate}
Then $K$ is isomorphic in $D(A)$ to a Koszul complex on a regular sequence
$f_1, \ldots, f_r \in A$ lifting the given elements
$\overline{f}_1, \ldots, \overline{f}_r$. Moreover, $A/(f_1, \ldots, f_r)$
is flat over $R$.
\end{lemma}

\begin{proof}
Let us use chain complexes in the proof of this lemma.
The Koszul complex $K_\bullet(\overline{f}_1, \ldots, \overline{f}_r)$
is defined in More on Algebra, Definition
\ref{more-algebra-definition-koszul-complex}.
By More on Algebra, Lemma \ref{more-algebra-lemma-lift-complex-stably-frees}
we can represent $K$ by a complex
$$
K_\bullet :
A \to A^{\oplus r} \to \ldots \to A^{\oplus r} \to A
$$
whose tensor product with $A/\mathfrak mA$ is equal (!)
to $K_\bullet(\overline{f}_1, \ldots, \overline{f}_r)$.
Denote $f_1, \ldots, f_r \in A$ the components of the
arrow $A^{\oplus r} \to A$. These $f_i$ are lifts of the
$\overline{f}_i$. By Algebra, Lemma
\ref{algebra-lemma-grothendieck-regular-sequence-general}
$f_1, \ldots, f_r$ form a regular sequence in $A$ and $A/(f_1, \ldots, f_r)$
is flat over $R$. Let $J = (f_1, \ldots, f_r) \subset A$.
Consider the diagram
$$
\xymatrix{
K_\bullet \ar[rd] \ar@{..>}[rr]_{\varphi_\bullet} & &
K_\bullet(f_1, \ldots, f_r) \ar[ld] \\
& A/J
}
$$
Since $f_1, \ldots, f_r$ is a regular sequence the south-west arrow
is a quasi-isomorphism (see
More on Algebra, Lemma \ref{more-algebra-lemma-regular-koszul-regular}).
Hence we can find the dotted arrow making the
diagram commute for example by
Algebra, Lemma \ref{algebra-lemma-compare-resolutions}.
Reducing modulo $\mathfrak m$ we obtain a commutative diagram
$$
\xymatrix{
K_\bullet(\overline{f}_1, \ldots, \overline{f}_r)
\ar[rd] \ar[rr]_{\overline{\varphi}_\bullet} & &
K_\bullet(\overline{f}_1, \ldots, \overline{f}_r) \ar[ld] \\
& (A/\mathfrak m A)/(\overline{f}_1, \ldots, \overline{f}_r)
}
$$
by our choice of $K_\bullet$. Thus $\overline{\varphi}$ is an isomorphism
in the derived category $D(A/\mathfrak m A)$. It follows that
$\overline{\varphi} \otimes_{A/\mathfrak m A}^\mathbf{L} \lambda$
is an isomorphism. Since $\overline{f}_i \in \mathfrak n / \mathfrak m A$
we see that
$$
\text{Tor}_i^{A/\mathfrak m A}(
K_\bullet(\overline{f}_1, \ldots, \overline{f}_r), \lambda)
=
K_i(\overline{f}_1, \ldots, \overline{f}_r) \otimes_{A/\mathfrak m A} \lambda
$$
Hence $\varphi_i \bmod \mathfrak n$ is invertible.
Since $A$ is local this means that $\varphi_i$ is an
isomorphism and the proof is complete.
\end{proof}

\end{document}
